\subsection{Non-Retentive Power-Gating Safety}
\label{sec:nonretentive-safety}

Power-gating the SA is non-retentive: values held in the processing-element
registers are not assumed to survive an OFF interval.  Consequently, an idle
interval alone is insufficient to authorize power gating.  The controller may
enter the OFF state only when it has both (i) evidence that the SA cannot be
needed before a certified deadline and (ii) evidence that losing the current
SA-resident state cannot change subsequent execution.

The second requirement follows from the weight-stationary instruction
semantics.  A \texttt{GEMM} instruction consumes weights already resident in
the SA and therefore cannot safely follow a non-retentive OFF interval.  A
\texttt{GEMM\_PRELOAD} instruction reloads the required weights before their
use.  Hence, for each of the two resident tile slots, the controller requires
the earliest future SA instruction to exist and to be
\texttt{GEMM\_PRELOAD}.  It also requires both resident slots to be occupied;
the future-tile bound described in Sec.~\ref{sec:causal-predictor}
limits the certificate to the earliest point at which a resident tile could
retire and previously unseen work could enter.

Table~\ref{tab:nonretentive-safety} summarizes the hazards checked by the
implemented controller.  These conditions are conjunctive: failure of any
entry condition suppresses a new OFF decision.

\begin{table}[t]
  \centering
  \caption{Safety conditions for a non-retentive SA OFF interval.}
  \label{tab:nonretentive-safety}
  \begin{tabular}{p{0.31\columnwidth}p{0.61\columnwidth}}
    \toprule
    Hazard & Implemented protection \\
    \midrule
    Useful SA work is already active
      & Require the SA to be inactive at every certificate stage and at
        entry commit. \\
    Committed work can reach the SA immediately
      & Require both the compute pipeline and execution-instruction queue to
        be empty. \\
    OFF destroys resident weights
      & Require the first future SA instruction in each resident tile to be
        \texttt{GEMM\_PRELOAD}, not \texttt{GEMM}. \\
    Unseen tile work arrives early
      & Require two resident tiles and cap the horizon using resident
        retirement and global work-drain information. \\
    Queue-distance evidence is ambiguous
      & Require exact memory-request tickets throughout certificate
        generation and at entry commit. \\
    Metadata changes after sampling
      & A write to either resident descriptor invalidates the certificate;
        a write during sleep forces wake-up. \\
    A saved deadline becomes stale
      & Freeze the committed deadline, decrement it monotonically, and never
        recompute or extend it while asleep. \\
    SA demand appears before the deadline
      & Wake immediately if the SA becomes active, either queue becomes
        nonempty, ticket exactness is lost, or the countdown reaches the
        one-cycle guard. \\
    \bottomrule
  \end{tabular}
\end{table}

Certificate construction is itself protected against stale observations.
The implementation uses three registered evaluation stages.  A descriptor
write, observed SA activity, a nonempty compute or execution queue, or loss of
ticket exactness invalidates the certificate at the intervening stage
boundaries.  The three elapsed evaluation cycles are deducted from the causal
lower bound.  With the evaluated 166-cycle break-even threshold and one-cycle
wake guard, a fresh bound must therefore be at least 170 cycles; after the
three-stage deduction, entry commits only when at least 167 cycles remain.

Once entry commits, the controller stores the certified remaining interval in
a countdown register.  While OFF, the controller cannot increase this value;
it either decrements the countdown or exits the OFF state.  This frozen-deadline
rule is important because continuously recomputing a later deadline from new
observations could incorrectly prolong an interval that was justified by an
older machine state.

We validate the decision logic by replaying each accepted architectural trace
against both the reference controller and its clock-gated implementation.  A
replay is accepted only when their emitted sleep episodes and summary records
match exactly.  Independently, the replay testbench counts a safety violation
whenever the controller requests sleep while the trace reports the SA active,
and terminates with an error if this count is nonzero.  Thus, the validation
establishes trace-level equivalence of the two controller implementations and
the absence of sleep--activity overlap in an accepted replay.

This validation has a deliberate scope.  The controller is evaluated in
shadow replay and does not feed a physical OFF/ON transition back into
ONNXim's scheduling or timing.  Therefore, unchanged workload cycle counts are
not, by themselves, a measurement of zero wake-up penalty in an integrated
power-gated NPU.  Device-level transition delay and energy are characterized
separately, and an end-to-end implementation that stalls execution during a
late wake-up would be required to measure any resulting performance penalty.
